% Standalone research entry 214; batch 208--217.
% Original section material preserved from Pasted text(3).txt.
% Research additions: 27 September 2026.
% Compile twice with: lualatex 214.tex
% !TeX program = lualatex
% Compile twice: lualatex open_problems_500.tex
% All references and prose are embedded; no BibTeX database or external inputs.
\documentclass[11pt,a4paper]{article}
\usepackage[margin=24mm,headheight=14pt]{geometry}
\usepackage{fontspec}
\setmainfont{Latin Modern Roman}
\setsansfont{Latin Modern Sans}
\setmonofont{Latin Modern Mono}
\usepackage{amsmath,amssymb,mathtools}
\usepackage{unicode-math}
\setmathfont{Latin Modern Math}
\usepackage{microtype}
\usepackage{enumitem,xurl,needspace}
\usepackage[dvipsnames]{xcolor}
\usepackage{hyperref}
\hypersetup{unicode=true,colorlinks=true,linkcolor=MidnightBlue,urlcolor=MidnightBlue,
 pdftitle={500 Mathematical Problem Statements},pdfauthor={Source-annotated compilation},bookmarksnumbered=true}
\usepackage{bookmark}
\usepackage{tocloft}
\setlength{\cftsecnumwidth}{3em}
\cftsetpnumwidth{2.5em}
\cftsetrmarg{4em}
\usepackage{titlesec}
\titleformat{\section}{\Large\bfseries\raggedright}{\thesection}{0.7em}{}
\usepackage{fancyhdr}
\pagestyle{fancy}\fancyhf{}
\fancyhead[L]{\small\sffamily Mathematical problem statements}
\fancyhead[R]{\small Source edition 22}
\fancyfoot[C]{\thepage}
\setlength{\parindent}{0pt}
\setlength{\parskip}{5pt plus 1pt}
\setlength{\emergencystretch}{3em}
\setcounter{tocdepth}{1}
\setcounter{secnumdepth}{1}
\setlist[itemize]{leftmargin=3em,itemsep=5pt,topsep=4pt,parsep=0pt}
\allowdisplaybreaks
\newcommand{\N}{\mathbb{N}}
\newcommand{\Z}{\mathbb{Z}}
\newcommand{\Q}{\mathbb{Q}}
\newcommand{\R}{\mathbb{R}}
\newcommand{\C}{\mathbb{C}}
\newcommand{\F}{\mathbb{F}}
\newcommand{\A}{\mathbb{A}}
\newcommand{\PP}{\mathbb P}
\newcommand{\E}{\mathbb{E}}
\newcommand{\eps}{\varepsilon}
\newcommand{\dd}{\,\mathrm d}
\newcommand{\sref}[2]{\hyperlink{source:#1:#2}{[S#2]}}
\newcommand{\eref}[2]{\hyperlink{extra:#1:#2}{[E#2]}}
% Turn the next switch off for a shorter reading copy. The quotations preserve
% every exactTarget field, including qualifications omitted from a short summary.
\newif\ifshowcatalogue
\showcataloguetrue
\newcommand{\cataloguescope}[1]{\ifshowcatalogue
 \paragraph{Exact scope in the supplied catalogue.}
 \begin{quote}\small #1\end{quote}\fi}

% Standalone wrapper; no external inputs or bibliography files are required.
\fancyhead[L]{\small\sffamily Research notebook: section 214}
\fancyhead[R]{\small 27 September 2026}
\hypersetup{pdftitle={Research attempt 214},pdfauthor={Research notebook}}
\setlength{\parskip}{4pt plus 1pt}
\setlist[itemize]{before=\small\interlinepenalty10000,itemsep=3pt}
\begin{document}
\setcounter{section}{213}
% ================================================================
% problemId: problem.global-existence-of-classical-solutions-to-the-relativistic-vlasov-maxwell-system
% source releaseRank: 214; source releaseStatus: open
% exactTarget SHA-256: dae7808465ab34a6042fe0611a67a250286da99311d5474474e3718b27a04a44
\clearpage
\section{\texorpdfstring{Global Existence of Classical Solutions to the Relativistic Vlasov-Maxwell System}{Global Existence of Classical Solutions to the Relativistic Vlasov-Maxwell System}}\label{214}
\subsection{Definitions and mathematical statement}
In units with particle mass and light speed one, set $\widehat v=v/\sqrt{1+|v|^2}$. A conventional one-species relativistic Vlasov--Maxwell system is
\[
 \partial_tf+\widehat v\cdot\nabla_xf+(E+\widehat v\times B)\cdot\nabla_vf=0,
 \quad \partial_tE=\nabla\times B-4\pi j,\quad\partial_tB=-\nabla\times E,
\]
with $\nabla\cdot E=4\pi\rho$, $\nabla\cdot B=0$, $\rho=\int f\dd v$, and $j=\int\widehat v f\dd v$. Finite energy means finiteness of
$\iint\sqrt{1+|v|^2}f\dd x\dd v+(8\pi)^{-1}\int(|E|^2+|B|^2)\dd x$.
The question is global classical evolution from all compatible data in the cited regularity class. The catalogue's phrase ``sufficiently regular'' does not fix one complete function-space hypothesis; this exposition supplies the equations but does not silently promote a particular data class to the exact target.
\par\smallskip\textit{Formulation and concept references:} \sref{214}{1}.
\cataloguescope{Prove or disprove that the Cauchy problem for the three-dimensional relativistic Vlasov-Maxwell system has a global classical solution (f,E,B) for every sufficiently regular finite-energy initial datum.}
\subsection{Short English statement}
Do all admissible sufficiently regular finite-energy data for the three-dimensional relativistic Vlasov–Maxwell system evolve as classical solutions for all time?
% BEGIN RESEARCH ADDITION 214
\subsection{Research attempt}

\paragraph{Current frontier.}
Luk--Strain's continuation theorem permits continuation if the projection
of the momentum support onto one fixed two-plane stays bounded; full
momentum-support boundedness is the earlier sufficient condition
\sref{214}{1}, Theorems 1.1--1.3. These are conditional regularity criteria,
not a priori estimates for arbitrary data. A 2026 two-paper program of
Xuecheng Wang treats arbitrarily large localized data with cylindrical
symmetry, with the second paper devoted to pointwise field estimates
\eref{214}{1}. That symmetry hypothesis cannot be removed by applying the
statement to general initial data. The conservation laws and null-cone
field decomposition in \sref{214}{1}, Sections 2--5, motivate the attack
below; its additional estimates are derived here.

\paragraph{Editorial scope note.}
The phrase ``sufficiently regular finite-energy initial datum'' in the
original statement does not specify a unique function-space class. I leave
it unchanged. To make every displayed integration and continuation step
precise, this attempt first works in the class explicitly used in
\sref{214}{1}: nonnegative compactly supported $f_0\in H^5(\R^6)$,
$E_0,B_0\in H^5(\R^3)$, with the stated divergence constraints. Proving
global existence in this subclass would not automatically treat arbitrary
noncompact velocity tails. The argument below does not close even for
unrestricted large data in this precise subclass.

\paragraph{Attack route.}
Three possibilities are examined: bound the electric work along particle
characteristics; propagate higher velocity moments and recover stronger
field norms; or exploit cancellation between nearly lightlike particles
and electromagnetic null components. The first would give continuation,
the second exposes exactly the integrability deficit, and the third tests
whether the deficit can be repaired by the geometry of the equations rather
than by an unsupported pointwise bound from energy.

\paragraph{Attempt.}
\textit{1. A precise sufficient electric-work estimate.}
Write $p_0(v)=\sqrt{1+|v|^2}$. Along a particle characteristic,
\[
 \dot X=\widehat V,\qquad \dot V=E(t,X)+\widehat V\times B(t,X),
 \qquad \frac{d}{dt}p_0(V)=\widehat V\cdot E(t,X).              \tag{214.1}
\]
The magnetic contribution cancels exactly; bounding the entire Lorentz force
can discard this useful structure. Let
\[
 P(t)=2+\sup\{p_0(v):f(s,x,v)\ne0,\ 0\le s\le t\}.
\]
Suppose for each finite $T$ there were constants $A_T,B_T<\infty$, depending
on the initial data and $T$ but not on an unknown bound for $P$, such that
all particle characteristics and all $t<T$ obeyed
\[
 \int_0^t(\widehat V(s)\cdot E(s,X(s)))_+\,ds
       \le A_T+B_T\log P(t).                                 \tag{214.2}
\]
Taking a supremum in (214.1) gives $P(t)\le C_T+B_T\log P(t)$.
Since $\log x\le\sqrt x$ for $x\ge1$, this implies the explicit bound
\[
 P(t)\le\left(\frac{B_T+\sqrt{B_T^2+4C_T}}2\right)^2.
\]
The compact-support continuation criterion would then extend the solution
past any finite maximal time \sref{214}{1}, Theorem 1.1. This proves the
implication from (214.2), not (214.2) itself. The positive part in that
estimate is a sufficient strengthening: cancellation in signed work might
allow a less restrictive successful approach.

\textit{2. Moment interpolation with the exponents retained.}
Put $F=\|f_0\|_\infty=\|f(t)\|_\infty$, and for $k\ge1$ define
\[
 m_k(t,x)=\int|v|^k f(t,x,v)\,dv,\qquad M_k(t)=\int m_k(t,x)\,dx.
\]
On any compact subinterval of the classical lifespan, integration by parts
in phase space is justified by the transported compact support. The
$v$-divergence of $E+\widehat v\times B$ vanishes, and
$v\cdot(\widehat v\times B)=0$. Consequently
\[
 M_k'(t)=k\iint |v|^{k-2}v\cdot E(t,x)f(t,x,v)\,dx\,dv.
                                                               \tag{214.3}
\]
For $k=1$ the formula follows by a smooth approximation of $|v|$ at zero.
Split the integral for $m_{k-1}$ at radius $R>0$:
\[
 m_{k-1}\le C_k F R^{k+2}+R^{-1}m_k.
\]
Optimizing $R$ (and handling $m_k=0$ by a limit) yields
\[
 m_{k-1}\le C_kF^{1/(k+3)}m_k^{(k+2)/(k+3)}.
\]
H\"older in $x$ with exponents $k+3$ and $(k+3)/(k+2)$ now gives
\[
 |M_k'|\le C_kF^{1/(k+3)}\|E(t)\|_{L^{k+3}_x}
                     M_k^{(k+2)/(k+3)}.
\]
Equivalently, by regularizing $M_k$ at zero and integrating,
\[
 M_k(t)^{1/(k+3)}\le M_k(0)^{1/(k+3)}+
 C_k F^{1/(k+3)}\int_0^t\|E(s)\|_{L^{k+3}_x}\,ds.             \tag{214.4}
\]
Conserved energy controls $E$ in $L^2_x$, not in $L^{k+3}_x$ for any
$k\ge1$. Even $k=1$ asks for $L^4_x$. Finite propagation in $x$ does not
reverse this inclusion on a bounded region. Equation (214.4) therefore
isolates a specific unproved field estimate rather than closing a moment
bootstrap. Also, finiteness of a finite moment is not itself boundedness of
the entire momentum support; a suitable moment continuation theorem would
still have to be invoked with its actual hypotheses.

\textit{3. Two concentration tests against an energy-only argument.}
Choose fixed smooth nonnegative compactly supported functions $\varphi,\psi$
with $\int\psi=1$, and define, for $R\ge2$,
\[
 f_{0,R}(x,v)=R^{-1}\varphi(x)\psi(v-Re_1).
\]
Its kinetic energy is bounded uniformly in $R$, its $L^\infty$ norm is
bounded, and its velocity support reaches size $R$. The density is
$R^{-1}\varphi$; take $B_{0,R}=0$ and the smooth Coulomb field
$E_{0,R}=R^{-1}\nabla\Delta^{-1}(4\pi\varphi)$. Then the constraints hold,
the field energy is $O(R^{-2})$, and each field belongs to $H^5$.
Thus total energy alone cannot bound the maximal initial momentum. This
example does not disprove a continuation bound depending on the actual
initial support.

A different test keeps spatial field energy bounded while the field on a
line becomes arbitrarily large. With smooth compactly supported
$\chi\in C_c^\infty(\R)$ and $\Psi\in C_c^\infty(\R^2)$, set
\[
 A_\delta(x)=(0,0,\chi(x_1)\Psi(x_2/\delta,x_3/\delta)),\qquad
 E_\delta=\nabla\times A_\delta.
\]
Then $\nabla\cdot E_\delta=0$ and
\[
 E_\delta=(\delta^{-1}\chi\,\partial_1\Psi,
                    -\chi'\Psi,0),\qquad
 \|E_\delta\|_2^2=C_1+\delta^2 C_2.
\]
Choose $\partial_1\Psi(0,0)\ne0$ and $\chi$ nonzero; the first component
on the $x_1$-axis has size comparable to $\delta^{-1}$. These are compatible
vacuum initial fields with $B_0=0$, and can also be added to a fixed Coulomb
field for a fixed particle density. They invalidate a universal
$L^\infty$ or line-trace bound deduced from $L^2$ energy alone. They do not
show large \emph{time-integrated} work along a characteristic of the evolved
coupled system: temporal oscillation and particle deflection are not
controlled by this initial-time construction.

\textit{4. The null cancellation and what it still leaves.}
For a unit vector $k$, put
$u=(|E|^2+|B|^2)/(8\pi)$ and $\mathcal S=(E\times B)/(4\pi)$.
Direct expansion gives the exact identity
\[
 u-k\cdot\mathcal S
 =\frac{|E\cdot k|^2+|B\cdot k|^2+|E_\perp+k\times B|^2}{8\pi}.
                                                               \tag{214.5}
\]
This is a nonnegative flux controlling particular field components. It is
the incoming-direction version of the null-cone flux in
\sref{214}{1}, Proposition 2.2: for its outward radial variable $\omega$,
use $k=-\omega$. This change is essential to keep the denominator
$1+\widehat v\cdot\omega$ consistent with the following $1-\widehat v\cdot k$.

Let $a=1-\widehat v\cdot k$. Then
\[
 1-|\widehat v|=\frac{1}{p_0(p_0+|v|)},\qquad
 |\widehat v-k|^2=2a-p_0^{-2}\le2a,
\]
and
\[
 E+\widehat v\times B=(E+k\times B)+(\widehat v-k)\times B.
                                                               \tag{214.6}
\]
The first term is controlled by the components in (214.5); the second gains
$\sqrt{2a}$ against an angular singularity, but still contains $B$.
This reproduces, in transparent algebraic form, why the most singular
parts of the field representation have better null components, while a
less singular part retains an uncontrolled magnetic factor; compare
\sref{214}{1}, Proposition 3.4.

At high momentum nearly parallel to $k$, $a$ can be of order $p_0^{-2}$.
The flux controls an $L^2$ integral on a cone, not a value on each timelike
particle curve. Applying Cauchy--Schwarz to the representation introduces
weighted angular particle-density integrals, precisely where control of the
projected momentum support is used in \sref{214}{1}, Sections 4--5. Neither
(214.5) nor total energy bounds those integrals in the required norm for
general data. Replacing the cone integral by a characteristic trace without
an additional estimate would therefore be the missing, unjustified step.

\paragraph{Stress test.}
All finite-time differentiations above are inside an existing classical
lifespan; no global solution is used to justify them. The concentration
examples are admissible data, not blow-up solutions. Their high Sobolev
norms may grow, so they do not refute estimates that legitimately depend on
those norms. The null identity is pointwise exact, but promoting it to the
uniform work bound (214.2) requires a nonlinear trace estimate not supplied
by energy conservation. A finite-dimensional angular projection cannot be
assumed bounded simply because the distribution is smooth at each time.
The regularity-class qualification and the distinction between weak and
classical solutions remain in force.

\paragraph{Outcome.}
\textbf{Outcome: Exploratory approach with unresolved gap.}

The attempt derives an explicit electric-work continuation implication, the
moment estimate (214.4) with its exact missing spatial exponent, and a
null-component cancellation consistent with the source's sign conventions.
Two smooth compatible concentration constructions eliminate energy-only
shortcuts without claiming a singularity. The remaining problem is a
self-consistent estimate of characteristic work or weighted angular density
which closes independently of the unknown momentum support. No global
large-data classical solution or counterexample is established.
% END RESEARCH ADDITION 214
\subsection{Sources}
\begin{itemize}\raggedright
\item[\textnormal{[S1]}]\hypertarget{source:214:1}{} arXiv:1406.0165, 'A new continuation criterion for the relativistic Vlasov-Maxwell system', 2014.
\url{https://arxiv.org/abs/1406.0165}.
{\small\textit{Catalogue formulation source.}}
% BEGIN RESEARCH REFERENCES 214
\item[\textnormal{[E1]}]\hypertarget{extra:214:1}{} Xuecheng Wang,
\emph{Large data global solution of the 3D RVM system with cylindrical symmetry I:
Iterative smoothing scheme}, arXiv:2203.01199, together with
\emph{II: Pointwise estimates}, arXiv:2607.14685 (16 July 2026), abstracts
and stated cylindrical-symmetry scope.
\url{https://arxiv.org/abs/2203.01199};
\url{https://arxiv.org/abs/2607.14685}.
{\small\textit{Recent large-data symmetry result; not a claim about arbitrary nonsymmetric data.}}
% END RESEARCH REFERENCES 214
\end{itemize}

\end{document}
